\section{The eleventh continuation: integral gates and distinct live targets}
\label{integral-status-section}

The inverse problem now has an exact arithmetic boundary. An actual
integer gcd divides the multiplier norm, but the norm-seven family
shows why this alone does not preserve a prescribed nonunit power
after primitive reduction. Opposite orientations give an exact local
reflection and a residual bill. On the power-free subdomain the
canonical output residual and root norm are determined explicitly.
The strict joint logarithmic threshold excludes cancellation, including
when the output root norm changes. These results identify a usable
integral domain; they do not produce compatible points in it.

For pure coefficients the positive rational locus has a complete
ordinary integral converse, with ramified and infinite inputs retained.
Together with the previous combined-square obstruction, this excludes
the even/even positive locus. The second quadratic-domain construction
provides simultaneous actual Gaussian and Eisenstein powers when the
first norm is a square. Its further biquadratic cover has eight branch
points. The even-exponent relation causes two geometric components;
the odd-exponent cover is geometrically connected. Neither its larger
genus nor its finite unit-twist list bounds the actual point heights
uniformly in the exponent. The odd branch remains a distinct live
arithmetic-geometric problem, with ramified and moving residual
branches also retained.

On the analytic route, the common large-cutoff set is independent of
the moving reciprocal caps. The finite signed estimate now transfers
to the global integer radical with explicit constants. A separate
prime-index theorem proves joint squarefreeness throughout a growing
prime interval beyond the root block. However, that interval's typical
full mass is sublinear; almost all full mass still lies in the farther
top-rank packet. Thus actual depth-one membership on the proved
window does not supply the missing signed bound on that packet.
Controlling deep excess together with its shallow-prime credit there,
and treating all exceptional roots, remains the primary analytic target.

The new joint formal inventory is 31 proved declarations: nine actual
content statements, seven residual-reflection arithmetic statements,
six integer-radical and logarithmic-threshold statements, and nine
global signed-bridge statements. The fresh run also compiles 53
unchanged local dependency declarations and audits all 84 axiom
queries against the pinned Lean and Mathlib environment. It reuses
the Mathlib compiled cache. One exact replay checks 24,192 coordinate
identities and 128 members of the norm-seven family. The other checks
the biquadratic polynomial identities, 119 finite relation kernels,
52 actual square-first seeds with complete elementary factorizations,
and 74 integer inverse cases. Those checks do not prove an infinite
or analytic theorem, or determine the full rational-point locus.

The complete oriented factorization, integral inverses, canonical
reflection, sieve estimates and geometric constructions have ordinary
proofs with independent internal review. They have not all been
formalized. In particular the formal local reflection premise is
explicit, and no uniform signed estimate or point-height bound is
introduced as an axiom. There is still no complete proof or disproof
of $abc$. Only the specific inverse assertions contradicted by the
actual examples are excluded; all unrefuted parent routes remain
available for further research.
